\section{Why This Document Is Inserted Here}

A080 treated units as bookkeeping: bridge constants, the separation of ratio
and absolute value. That is the \emph{operative} side. This document provides
the \emph{fundamental} side that A080 presupposes but does not develop: why a
natural unit system exists in this theory at all, in which all dimensionful
constants take the value one.

It precedes the gravitation document (A145) and quantum mechanics (A160)
because both use it: $G$ is obtained from $\xi$ and a mass, and that is a
statement about structure rather than numerical conventions only in natural units.

\section{The Origin Lies in the Duality}

\textbf{[K]} The natural unit system of the theory is not a choice but a
consequence of the first stipulation. From $\tilde T\cdot m = 1$ (A020) it
follows that time and mass are reciprocal -- the same quantity in two readings.
The conversion between them is therefore not a physical constant but the
identity.

\textbf{[B]} In standard physics three constants perform the conversion between
four fundamental dimensions:
\[
  c:\ \text{length}\leftrightarrow\text{time},\quad
  \hbar:\ \text{mass}\leftrightarrow\text{time},\quad
  k_B:\ \text{temperature}\leftrightarrow\text{energy}.
\]
But once time and mass are the same quantity, the system collapses: there is
only one independent dimension left, and these three conversion factors become
one.

\textbf{[STIPULATION]} A word of demarcation that belongs here: this mechanism
concerns the \emph{dimensionful} bridge constants $c, \hbar, k_B$ (and, via
the mass-length connection, $G$; A145). The fine-structure constant $\alpha$
is \emph{not} among them. It is dimensionless, and its reduction to one has a
different reason -- electromagnetic duality, not time-mass duality. That is a
separate chapter and is treated in A130, not here. This document claims nothing
about $\alpha$; it says only why $c, \hbar, k_B$ become one.

\textbf{[K]} In natural units therefore
\[
  [M] = [E] = [T^{-1}] = [L^{-1}],
\]
mass, energy, inverse time, and inverse length are a single dimension. This is
not set simplifyingly but is the direct consequence of the duality.

\section{Why This Is More Than a Convention}

\textbf{[B]} One can set $\hbar = c = 1$ in any theory; that alone is
inconsequential. The difference here is that the \emph{reduction to a single
dimension} is a substantive claim, not a notation.

In standard physics, after $\hbar = c = 1$ one mass dimension remains -- energy
is a genuine unit that must be chosen. Here \emph{nothing} remains to choose:
the one remaining dimension is scaled by $\xi$ (A020), and the absolute scale
is the only remaining input (A080, A130).

\textbf{[STIPULATION]} The bookkeeping of A080 is thus bound to its ground:
the bridge constants are not equal to one because one sets them so, but because
the duality identifies their dimensions. Setting $=1$ is reading off an identity,
not a normalisation.

\section{The Four Dresses, Read Fundamentally}

\textbf{[K]} A080 introduced the identity
\[
  m_0 c^2 = E_0 = k_B T_0 = \hbar/\tilde T
\]
as an exact conversion. Here is its reason: the four expressions are not four
quantities that happen to be equal, but one quantity in four dresses, because
the duality places mass, energy, temperature, and inverse time on the same
dimension.

The conversion factors $c^2$, $k_B$, $\hbar$ are the seams between the dresses.
In natural units the seams disappear, because there is only one garment.

\section{What Follows for the Following Documents}

\textbf{[B]} Three consequences that A145 and A160 presuppose:

First, every mass is an inverse length and vice versa. This allows a
gravitational coupling linking masses with lengths to be expressed through $\xi$
alone -- the content of A145.

Second, every dimensionful constant is either one or a calibration. There is no
third case. Where in the corpus such a constant carries another value, that is
an SI artefact, not a fact of nature.

Third, ratios are invariant under the choice of scale (A080), and the choice
of scale is the \emph{only} remaining degree of freedom in the unit system.
Everything dimensionless is thereby scale-free-determined or not at all.

\section{Verification Script}

\begin{itemize}
  \item That $\tilde T\cdot m=1$ reduces the four-dimension system to one
    dimension, shown in the dimensional bookkeeping; and that the Planck
    quantities are exactly the conversion factors of the reduced system:
    \texttt{python/A\_Serie\_Skripte/a085\_natuerliche\_einheiten.py}
\end{itemize}

\section{Open Questions}

\textbf{First, the absolute scale remains an input.} The reduction to one
dimension is derived; what size this one dimension has is not. This is the same
gap as in A020 (order of magnitude of $\xi$) and A130 (calibration of $E_0$),
here in its most general form.

\textbf{Second, the identification of temperature with energy is an additional
thermodynamic assumption.} $k_B=1$ presupposes that temperature is an energy
per degree of freedom; that is standard, but it does not follow from the
duality, it stands beside it.

\section{Supersedes}

This document supplements A080 with the fundamental justification and subsumes
the corresponding parts of: Dok.~013 (SI), Dok.~014 (natural/SI), Dok.~015
(systematics of natural units), Dok.~012 (units section). Incorporated: P40.

\section{Use of AI Tools}

This document was created with the assistance of AI tools.
Responsibility for content and errors rests entirely with the author.
Full wording of the rules in A010.
